% ============================================================================
% Глава 1: Въведение
% ============================================================================

\section{ВЪВЕДЕНИЕ}

% TODO: В тази глава трябва да представите проекта и да обясните контекста, в който той се развива.
% TODO: Какъв е бизнес проблемът, който разрешавате? - Опишете конкретен проблем, свързан със съхранението, управлението или обработката на данни, който възниква в дадена индустрия или организация. Например, ако разработвате система за електронна търговия, може да имате проблем със съхраняването на динамични продуктови каталози. Ако работите върху IoT платформа, може да се сблъскате с нуждата от обработка на големо количество сензорни данни в реално време.
% TODO: Какви са целите на проекта? - Дефинирайте какво очаквате да постигнете с решението си. Това може да включва подобряване на скоростта на заявките, увеличаване на мащабируемостта или улесняване на достъпа до данните.
% TODO: Защо избрахте нерелационна база данни? - Обяснете защо NoSQL моделът е по-подходящ за вашия проект в сравнение с релационните бази данни. Опишете ключовите предимства, като гъвкавост в структурирането на данните, по-добра производителност при работа с големи обеми информация или възможности за хоризонтално мащабиране.